% Options for packages loaded elsewhere
% Options for packages loaded elsewhere
\PassOptionsToPackage{unicode}{hyperref}
\PassOptionsToPackage{hyphens}{url}
\PassOptionsToPackage{dvipsnames,svgnames,x11names}{xcolor}
%
\documentclass[
  letterpaper,
  DIV=11,
  numbers=noendperiod]{scrartcl}
\usepackage{xcolor}
\usepackage[margin=2.5cm]{geometry}
\usepackage{amsmath,amssymb}
\setcounter{secnumdepth}{5}
\usepackage{iftex}
\ifPDFTeX
  \usepackage[T1]{fontenc}
  \usepackage[utf8]{inputenc}
  \usepackage{textcomp} % provide euro and other symbols
\else % if luatex or xetex
  \usepackage{unicode-math} % this also loads fontspec
  \defaultfontfeatures{Scale=MatchLowercase}
  \defaultfontfeatures[\rmfamily]{Ligatures=TeX,Scale=1}
\fi
\usepackage{lmodern}
\ifPDFTeX\else
  % xetex/luatex font selection
\fi
% Use upquote if available, for straight quotes in verbatim environments
\IfFileExists{upquote.sty}{\usepackage{upquote}}{}
\IfFileExists{microtype.sty}{% use microtype if available
  \usepackage[]{microtype}
  \UseMicrotypeSet[protrusion]{basicmath} % disable protrusion for tt fonts
}{}
\makeatletter
\@ifundefined{KOMAClassName}{% if non-KOMA class
  \IfFileExists{parskip.sty}{%
    \usepackage{parskip}
  }{% else
    \setlength{\parindent}{0pt}
    \setlength{\parskip}{6pt plus 2pt minus 1pt}}
}{% if KOMA class
  \KOMAoptions{parskip=half}}
\makeatother
% Make \paragraph and \subparagraph free-standing
\makeatletter
\ifx\paragraph\undefined\else
  \let\oldparagraph\paragraph
  \renewcommand{\paragraph}{
    \@ifstar
      \xxxParagraphStar
      \xxxParagraphNoStar
  }
  \newcommand{\xxxParagraphStar}[1]{\oldparagraph*{#1}\mbox{}}
  \newcommand{\xxxParagraphNoStar}[1]{\oldparagraph{#1}\mbox{}}
\fi
\ifx\subparagraph\undefined\else
  \let\oldsubparagraph\subparagraph
  \renewcommand{\subparagraph}{
    \@ifstar
      \xxxSubParagraphStar
      \xxxSubParagraphNoStar
  }
  \newcommand{\xxxSubParagraphStar}[1]{\oldsubparagraph*{#1}\mbox{}}
  \newcommand{\xxxSubParagraphNoStar}[1]{\oldsubparagraph{#1}\mbox{}}
\fi
\makeatother

\usepackage{color}
\usepackage{fancyvrb}
\newcommand{\VerbBar}{|}
\newcommand{\VERB}{\Verb[commandchars=\\\{\}]}
\DefineVerbatimEnvironment{Highlighting}{Verbatim}{commandchars=\\\{\}}
% Add ',fontsize=\small' for more characters per line
\usepackage{framed}
\definecolor{shadecolor}{RGB}{241,243,245}
\newenvironment{Shaded}{\begin{snugshade}}{\end{snugshade}}
\newcommand{\AlertTok}[1]{\textcolor[rgb]{0.68,0.00,0.00}{#1}}
\newcommand{\AnnotationTok}[1]{\textcolor[rgb]{0.37,0.37,0.37}{#1}}
\newcommand{\AttributeTok}[1]{\textcolor[rgb]{0.40,0.45,0.13}{#1}}
\newcommand{\BaseNTok}[1]{\textcolor[rgb]{0.68,0.00,0.00}{#1}}
\newcommand{\BuiltInTok}[1]{\textcolor[rgb]{0.00,0.23,0.31}{#1}}
\newcommand{\CharTok}[1]{\textcolor[rgb]{0.13,0.47,0.30}{#1}}
\newcommand{\CommentTok}[1]{\textcolor[rgb]{0.37,0.37,0.37}{#1}}
\newcommand{\CommentVarTok}[1]{\textcolor[rgb]{0.37,0.37,0.37}{\textit{#1}}}
\newcommand{\ConstantTok}[1]{\textcolor[rgb]{0.56,0.35,0.01}{#1}}
\newcommand{\ControlFlowTok}[1]{\textcolor[rgb]{0.00,0.23,0.31}{\textbf{#1}}}
\newcommand{\DataTypeTok}[1]{\textcolor[rgb]{0.68,0.00,0.00}{#1}}
\newcommand{\DecValTok}[1]{\textcolor[rgb]{0.68,0.00,0.00}{#1}}
\newcommand{\DocumentationTok}[1]{\textcolor[rgb]{0.37,0.37,0.37}{\textit{#1}}}
\newcommand{\ErrorTok}[1]{\textcolor[rgb]{0.68,0.00,0.00}{#1}}
\newcommand{\ExtensionTok}[1]{\textcolor[rgb]{0.00,0.23,0.31}{#1}}
\newcommand{\FloatTok}[1]{\textcolor[rgb]{0.68,0.00,0.00}{#1}}
\newcommand{\FunctionTok}[1]{\textcolor[rgb]{0.28,0.35,0.67}{#1}}
\newcommand{\ImportTok}[1]{\textcolor[rgb]{0.00,0.46,0.62}{#1}}
\newcommand{\InformationTok}[1]{\textcolor[rgb]{0.37,0.37,0.37}{#1}}
\newcommand{\KeywordTok}[1]{\textcolor[rgb]{0.00,0.23,0.31}{\textbf{#1}}}
\newcommand{\NormalTok}[1]{\textcolor[rgb]{0.00,0.23,0.31}{#1}}
\newcommand{\OperatorTok}[1]{\textcolor[rgb]{0.37,0.37,0.37}{#1}}
\newcommand{\OtherTok}[1]{\textcolor[rgb]{0.00,0.23,0.31}{#1}}
\newcommand{\PreprocessorTok}[1]{\textcolor[rgb]{0.68,0.00,0.00}{#1}}
\newcommand{\RegionMarkerTok}[1]{\textcolor[rgb]{0.00,0.23,0.31}{#1}}
\newcommand{\SpecialCharTok}[1]{\textcolor[rgb]{0.37,0.37,0.37}{#1}}
\newcommand{\SpecialStringTok}[1]{\textcolor[rgb]{0.13,0.47,0.30}{#1}}
\newcommand{\StringTok}[1]{\textcolor[rgb]{0.13,0.47,0.30}{#1}}
\newcommand{\VariableTok}[1]{\textcolor[rgb]{0.07,0.07,0.07}{#1}}
\newcommand{\VerbatimStringTok}[1]{\textcolor[rgb]{0.13,0.47,0.30}{#1}}
\newcommand{\WarningTok}[1]{\textcolor[rgb]{0.37,0.37,0.37}{\textit{#1}}}

\usepackage{longtable,booktabs,array}
\usepackage{calc} % for calculating minipage widths
% Correct order of tables after \paragraph or \subparagraph
\usepackage{etoolbox}
\makeatletter
\patchcmd\longtable{\par}{\if@noskipsec\mbox{}\fi\par}{}{}
\makeatother
% Allow footnotes in longtable head/foot
\IfFileExists{footnotehyper.sty}{\usepackage{footnotehyper}}{\usepackage{footnote}}
\makesavenoteenv{longtable}
\usepackage{graphicx}
\makeatletter
\newsavebox\pandoc@box
\newcommand*\pandocbounded[1]{% scales image to fit in text height/width
  \sbox\pandoc@box{#1}%
  \Gscale@div\@tempa{\textheight}{\dimexpr\ht\pandoc@box+\dp\pandoc@box\relax}%
  \Gscale@div\@tempb{\linewidth}{\wd\pandoc@box}%
  \ifdim\@tempb\p@<\@tempa\p@\let\@tempa\@tempb\fi% select the smaller of both
  \ifdim\@tempa\p@<\p@\scalebox{\@tempa}{\usebox\pandoc@box}%
  \else\usebox{\pandoc@box}%
  \fi%
}
% Set default figure placement to htbp
\def\fps@figure{htbp}
\makeatother





\setlength{\emergencystretch}{3em} % prevent overfull lines

\providecommand{\tightlist}{%
  \setlength{\itemsep}{0pt}\setlength{\parskip}{0pt}}



 


\usepackage{booktabs}
\usepackage{longtable}
\usepackage{array}
\usepackage{multirow}
\usepackage{wrapfig}
\usepackage{float}
\usepackage{colortbl}
\usepackage{pdflscape}
\usepackage{tabu}
\usepackage{threeparttable}
\usepackage{threeparttablex}
\usepackage[normalem]{ulem}
\usepackage{makecell}
\usepackage{xcolor}
\usepackage{caption}
\usepackage{anyfontsize}
\KOMAoption{captions}{tableheading}
\usepackage{booktabs}
\makeatletter
\@ifpackageloaded{caption}{}{\usepackage{caption}}
\AtBeginDocument{%
\ifdefined\contentsname
  \renewcommand*\contentsname{Table of contents}
\else
  \newcommand\contentsname{Table of contents}
\fi
\ifdefined\listfigurename
  \renewcommand*\listfigurename{List of Figures}
\else
  \newcommand\listfigurename{List of Figures}
\fi
\ifdefined\listtablename
  \renewcommand*\listtablename{List of Tables}
\else
  \newcommand\listtablename{List of Tables}
\fi
\ifdefined\figurename
  \renewcommand*\figurename{Figure}
\else
  \newcommand\figurename{Figure}
\fi
\ifdefined\tablename
  \renewcommand*\tablename{Table}
\else
  \newcommand\tablename{Table}
\fi
}
\@ifpackageloaded{float}{}{\usepackage{float}}
\floatstyle{ruled}
\@ifundefined{c@chapter}{\newfloat{codelisting}{h}{lop}}{\newfloat{codelisting}{h}{lop}[chapter]}
\floatname{codelisting}{Listing}
\newcommand*\listoflistings{\listof{codelisting}{List of Listings}}
\makeatother
\makeatletter
\makeatother
\makeatletter
\@ifpackageloaded{caption}{}{\usepackage{caption}}
\@ifpackageloaded{subcaption}{}{\usepackage{subcaption}}
\makeatother
\usepackage{bookmark}
\IfFileExists{xurl.sty}{\usepackage{xurl}}{} % add URL line breaks if available
\urlstyle{same}
\hypersetup{
  pdftitle={Analyses de l'essai Modafinil (Données Fictives)},
  pdfauthor={Thomas Husson - Groupe 52},
  colorlinks=true,
  linkcolor={blue},
  filecolor={Maroon},
  citecolor={Blue},
  urlcolor={Blue},
  pdfcreator={LaTeX via pandoc}}


\title{Analyses de l'essai Modafinil (Données Fictives)}
\author{Thomas Husson - Groupe 52}
\date{}
\begin{document}
\maketitle

\renewcommand*\contentsname{Table of contents}
{
\hypersetup{linkcolor=}
\setcounter{tocdepth}{3}
\tableofcontents
}

\section{Introduction}\label{introduction}

Ce document présente les résultats statistiques basés sur les données
fictives pour l'essai croisé comparant le Modafinil et un Placebo chez
des patients atteints du syndrome post-polio.

Le jeu de données fictif de 38 patients a été utilisé pour illustrer la
méthodologie proposée.

\section{Puissance statistique}\label{puissance-statistique}

Voici le calcul de puissance effectué sur le principe d'un essai croisé
:

\begin{Shaded}
\begin{Highlighting}[]
\FunctionTok{pwr.t.test}\NormalTok{(}\AttributeTok{d =} \DecValTok{1}\SpecialCharTok{/}\DecValTok{1}\NormalTok{, }\AttributeTok{sig.level =} \FloatTok{0.05}\NormalTok{, }\AttributeTok{power =} \FloatTok{0.80}\NormalTok{, }\AttributeTok{type =} \StringTok{"paired"}\NormalTok{)}
\end{Highlighting}
\end{Shaded}

\begin{verbatim}

     Paired t test power calculation 

              n = 9.93785
              d = 1
      sig.level = 0.05
          power = 0.8
    alternative = two.sided

NOTE: n is number of *pairs*
\end{verbatim}

\textbf{Interprétation détaillée de l'output de l'analyse de puissance
:} L'output de la fonction \texttt{pwr.t.test} fournit les éléments de
dimensionnement suivants : - \textbf{\texttt{n\ =\ 9.936}} : Dans le
contexte d'un modèle de test de comparaison apparié
(\texttt{type\ =\ "paired"}), cette valeur \texttt{n} correspond de fait
au nombre de \emph{paires} nécessaires. Il faudrait inclure un seuil
minimal de 10 patients (qui reçoivent tous les deux traitements dans le
cadre du plan croisé) pour atteindre l'objectif. -
\textbf{\texttt{d\ =\ 1}} : Il s'agit de la taille de l'effet
standardisée stipulée (Taille d'effet de Cohen). Un \(d=1\) correspond
cliniquement à un effet particulièrement ``fort'' (une différence
moyenne d'exactement 1 écart type entier séparant les deux groupes). -
\textbf{\texttt{power\ =\ 0.8}} : Cette valeur a été fixée a priori; il
s'agit de la puissance probabiliste statistique d'écarter l'hypothèse
nulle \(H_0\) (ne pas rater une vraie différence). -
\textbf{\texttt{sig.level\ =\ 0.05}} : Le test s'effectue avec un risque
alpha d'erreur de première espèce de 5\%. - \textbf{Conclusion globale
:} Notre base de données fictive comportant un effectif de \(N = 38\)
individus pour une exigence de base de 10 patients, l'étude s'avère
manifestement expérimentalement sur-dimensionnée pour des effets forts.
Sa très bonne puissance assurera l'absence d'erreurs de type 2.

\section{Préparation et description des
données}\label{pruxe9paration-et-description-des-donnuxe9es}

\begin{Shaded}
\begin{Highlighting}[]
\NormalTok{csv\_path }\OtherTok{\textless{}{-}} \StringTok{"df\_large\_fictif\_38.csv"}
\NormalTok{df\_large\_simplifie }\OtherTok{\textless{}{-}} \FunctionTok{read.csv}\NormalTok{(csv\_path, }\AttributeTok{stringsAsFactors =} \ConstantTok{FALSE}\NormalTok{) }\SpecialCharTok{\%\textgreater{}\%}
    \FunctionTok{mutate}\NormalTok{(}
        \AttributeTok{centre =} \FunctionTok{factor}\NormalTok{(centre, }\AttributeTok{levels =} \FunctionTok{c}\NormalTok{(}\StringTok{"NIH"}\NormalTok{, }\StringTok{"NRH"}\NormalTok{, }\StringTok{"WRAMC"}\NormalTok{)),}
        \AttributeTok{sequence\_randomisation =} \FunctionTok{factor}\NormalTok{(sequence\_randomisation, }\AttributeTok{levels =} \FunctionTok{c}\NormalTok{(}\StringTok{"MP"}\NormalTok{, }\StringTok{"PM"}\NormalTok{)),}
        \AttributeTok{sexe =} \FunctionTok{factor}\NormalTok{(sexe, }\AttributeTok{levels =} \FunctionTok{c}\NormalTok{(}\StringTok{"Homme"}\NormalTok{, }\StringTok{"Femme"}\NormalTok{))}
\NormalTok{    )}

\NormalTok{id\_vars }\OtherTok{\textless{}{-}} \FunctionTok{names}\NormalTok{(df\_large\_simplifie)[}\SpecialCharTok{!}\FunctionTok{grepl}\NormalTok{(}\StringTok{"\_(S0|S3|S6|S8|S11|S14)$"}\NormalTok{, }\FunctionTok{names}\NormalTok{(df\_large\_simplifie))]}

\NormalTok{df\_long\_simplifie }\OtherTok{\textless{}{-}}\NormalTok{ reshape2}\SpecialCharTok{::}\FunctionTok{melt}\NormalTok{(}
\NormalTok{    df\_large\_simplifie,}
    \AttributeTok{id.vars =}\NormalTok{ id\_vars,}
    \AttributeTok{variable.name =} \StringTok{"variable\_visite"}\NormalTok{,}
    \AttributeTok{value.name =} \StringTok{"valeur"}
\NormalTok{) }\SpecialCharTok{\%\textgreater{}\%}
    \FunctionTok{mutate}\NormalTok{(}
        \AttributeTok{variable =} \FunctionTok{sub}\NormalTok{(}\StringTok{"\_(S0|S3|S6|S8|S11|S14)$"}\NormalTok{, }\StringTok{""}\NormalTok{, variable\_visite),}
        \AttributeTok{visite =} \FunctionTok{sub}\NormalTok{(}\StringTok{"\^{}.*\_(S0|S3|S6|S8|S11|S14)$"}\NormalTok{, }\StringTok{"}\SpecialCharTok{\textbackslash{}\textbackslash{}}\StringTok{1"}\NormalTok{, variable\_visite)}
\NormalTok{    ) }\SpecialCharTok{\%\textgreater{}\%}
\NormalTok{    reshape2}\SpecialCharTok{::}\FunctionTok{dcast}\NormalTok{(}
        \AttributeTok{formula =} \FunctionTok{as.formula}\NormalTok{(}\FunctionTok{paste}\NormalTok{(}\FunctionTok{paste}\NormalTok{(}\FunctionTok{c}\NormalTok{(id\_vars, }\StringTok{"visite"}\NormalTok{), }\AttributeTok{collapse =} \StringTok{" + "}\NormalTok{), }\StringTok{"\textasciitilde{} variable"}\NormalTok{)),}
        \AttributeTok{value.var =} \StringTok{"valeur"}\NormalTok{,}
        \AttributeTok{fun.aggregate =} \ControlFlowTok{function}\NormalTok{(x) x[}\DecValTok{1}\NormalTok{]}
\NormalTok{    ) }\SpecialCharTok{\%\textgreater{}\%}
    \FunctionTok{mutate}\NormalTok{(}
        \AttributeTok{visite =} \FunctionTok{factor}\NormalTok{(visite, }\AttributeTok{levels =} \FunctionTok{c}\NormalTok{(}\StringTok{"S0"}\NormalTok{, }\StringTok{"S3"}\NormalTok{, }\StringTok{"S6"}\NormalTok{, }\StringTok{"S8"}\NormalTok{, }\StringTok{"S11"}\NormalTok{, }\StringTok{"S14"}\NormalTok{), }\AttributeTok{ordered =} \ConstantTok{TRUE}\NormalTok{),}
        \AttributeTok{periode =} \FunctionTok{case\_when}\NormalTok{(visite }\SpecialCharTok{\%in\%} \FunctionTok{c}\NormalTok{(}\StringTok{"S0"}\NormalTok{, }\StringTok{"S3"}\NormalTok{, }\StringTok{"S6"}\NormalTok{) }\SpecialCharTok{\textasciitilde{}} \DecValTok{1}\NormalTok{L, visite }\SpecialCharTok{\%in\%} \FunctionTok{c}\NormalTok{(}\StringTok{"S8"}\NormalTok{, }\StringTok{"S11"}\NormalTok{, }\StringTok{"S14"}\NormalTok{) }\SpecialCharTok{\textasciitilde{}} \DecValTok{2}\NormalTok{L),}
        \AttributeTok{traitement =} \FunctionTok{case\_when}\NormalTok{(}
\NormalTok{            sequence\_randomisation }\SpecialCharTok{==} \StringTok{"MP"} \SpecialCharTok{\&}\NormalTok{ periode }\SpecialCharTok{==} \DecValTok{1}\NormalTok{L }\SpecialCharTok{\textasciitilde{}} \StringTok{"Modafinil"}\NormalTok{,}
\NormalTok{            sequence\_randomisation }\SpecialCharTok{==} \StringTok{"MP"} \SpecialCharTok{\&}\NormalTok{ periode }\SpecialCharTok{==} \DecValTok{2}\NormalTok{L }\SpecialCharTok{\textasciitilde{}} \StringTok{"Placebo"}\NormalTok{,}
\NormalTok{            sequence\_randomisation }\SpecialCharTok{==} \StringTok{"PM"} \SpecialCharTok{\&}\NormalTok{ periode }\SpecialCharTok{==} \DecValTok{1}\NormalTok{L }\SpecialCharTok{\textasciitilde{}} \StringTok{"Placebo"}\NormalTok{,}
\NormalTok{            sequence\_randomisation }\SpecialCharTok{==} \StringTok{"PM"} \SpecialCharTok{\&}\NormalTok{ periode }\SpecialCharTok{==} \DecValTok{2}\NormalTok{L }\SpecialCharTok{\textasciitilde{}} \StringTok{"Modafinil"}
\NormalTok{        ),}
        \AttributeTok{traitement =} \FunctionTok{factor}\NormalTok{(traitement, }\AttributeTok{levels =} \FunctionTok{c}\NormalTok{(}\StringTok{"Placebo"}\NormalTok{, }\StringTok{"Modafinil"}\NormalTok{))}
\NormalTok{    )}

\NormalTok{cols\_num }\OtherTok{\textless{}{-}} \FunctionTok{c}\NormalTok{(}\StringTok{"fss"}\NormalTok{, }\StringTok{"vasf"}\NormalTok{, }\StringTok{"fis"}\NormalTok{, }\StringTok{"sf36\_pcs"}\NormalTok{, }\StringTok{"sf36\_mcs"}\NormalTok{)}
\NormalTok{df\_long\_simplifie[cols\_num] }\OtherTok{\textless{}{-}} \FunctionTok{lapply}\NormalTok{(df\_long\_simplifie[cols\_num], as.numeric)}
\end{Highlighting}
\end{Shaded}

\subsection{Caractéristiques initiales des
patients}\label{caractuxe9ristiques-initiales-des-patients}

\begin{table}
\fontsize{12.0pt}{14.0pt}\selectfont
\begin{tabular*}{\linewidth}{@{\extracolsep{\fill}}lcccc}
\toprule
\textbf{Characteristic} & \textbf{Overall}  N = 38\textsuperscript{\textit{1}} & \textbf{MP}  N = 19\textsuperscript{\textit{1}} & \textbf{PM}  N = 19\textsuperscript{\textit{1}} & \textbf{p-value}\textsuperscript{\textit{2}} \\ 
\midrule\addlinespace[2.5pt]
Âge (années) & 65 (50, 70) & 64 (53, 72) & 65 (46, 69) & 0.5 \\ 
Sexe &  &  &  & 0.7 \\ 
    Homme & 23 (61\%) & 11 (58\%) & 12 (63\%) &  \\ 
    Femme & 15 (39\%) & 8 (42\%) & 7 (37\%) &  \\ 
Centre d'inclusion &  &  &  & >0.9 \\ 
    NIH & 13 (34\%) & 7 (37\%) & 6 (32\%) &  \\ 
    NRH & 13 (34\%) & 6 (32\%) & 7 (37\%) &  \\ 
    WRAMC & 12 (32\%) & 6 (32\%) & 6 (32\%) &  \\ 
\bottomrule
\end{tabular*}
\begin{minipage}{\linewidth}
\vspace{.05em}
\textsuperscript{\textit{1}}Median (Q1, Q3); n (\%)\\
\textsuperscript{\textit{2}}Wilcoxon rank sum test; Pearson's Chi-squared test\\
\end{minipage}
\end{table}

\textbf{Interprétation détaillée du tableau généré par
\texttt{gtsummary} :} L'output fournit l'état des lieux biométriques
initial de notre échantillon stratifié selon la
\texttt{séquence\ de\ randomisation}. - \textbf{Statistiques
descriptives (``PM'' et ``MP'') :} Le tableau reporte les effectifs
(N=19 pour chaque groupe, prouvant la balance du tirage), la médiane
avec l'écart inter-quartile {[}IQR{]} (ex: âge médian de 50-51 ans) pour
les variables continues, et la proportion pour les variables nominales.
- \textbf{La colonne ``p-value'' (\texttt{add\_p()}) :} Elle teste de
façon non paramétrique (Wilcoxon pour l'âge, exact de Fisher pour sexe
et centres) l'hypothèse nulle \(H_0\) de parfaite indépendance entre le
caractère individuel et sa séquence d'inclusion. - \textbf{Conclusion
globale :} Puisqu'aucune de ces p-values n'est inférieure au seuil
strict \(\alpha\) de 0.05, on en tire le postulat crucial affirmant
l'absence de biais. Les groupes de départ sont considérés
fondamentalement identiques pour l'induction. Le protocole de
randomisation a fonctionné avec succès.

\section{Analyses du critère principal : FSS (Fatigue Severity
Scale)}\label{analyses-du-crituxe8re-principal-fss-fatigue-severity-scale}

L'analyse de l'effet du traitement sur la fatigue est réalisée à l'aide
d'un modèle mixte tenant compte de l'effet aléatoire du patient, ainsi
que des effets fixes période, traitement, et séquence de randomisation.
Cette analyse s'effectue sur les données des visites en temps S6 et S14
(fin de chaque période de traitement de 6 semaines).

\begin{Shaded}
\begin{Highlighting}[]
\NormalTok{modele\_principal }\OtherTok{\textless{}{-}}\NormalTok{ lme4}\SpecialCharTok{::}\FunctionTok{lmer}\NormalTok{(}
\NormalTok{  fss }\SpecialCharTok{\textasciitilde{}}\NormalTok{ traitement }\SpecialCharTok{+}\NormalTok{ periode }\SpecialCharTok{+}\NormalTok{ sequence\_randomisation }\SpecialCharTok{+}\NormalTok{ (}\DecValTok{1} \SpecialCharTok{|}\NormalTok{ id\_patient),}
  \AttributeTok{data =} \FunctionTok{subset}\NormalTok{(df\_long\_simplifie, visite }\SpecialCharTok{\%in\%} \FunctionTok{c}\NormalTok{(}\StringTok{"S6"}\NormalTok{, }\StringTok{"S14"}\NormalTok{))}
\NormalTok{)}
\FunctionTok{summary}\NormalTok{(modele\_principal)}
\end{Highlighting}
\end{Shaded}

\begin{verbatim}
Linear mixed model fit by REML ['lmerMod']
Formula: fss ~ traitement + periode + sequence_randomisation + (1 | id_patient)
   Data: subset(df_long_simplifie, visite %in% c("S6", "S14"))

REML criterion at convergence: 115.9

Scaled residuals: 
     Min       1Q   Median       3Q      Max 
-1.45580 -0.51285 -0.02354  0.55839  1.87282 

Random effects:
 Groups     Name        Variance Std.Dev.
 id_patient (Intercept) 0.58197  0.7629  
 Residual               0.04901  0.2214  
Number of obs: 76, groups:  id_patient, 38

Fixed effects:
                         Estimate Std. Error t value
(Intercept)               4.97368    0.19588  25.392
traitementModafinil      -0.98947    0.05079 -19.483
periode                   0.14211    0.05079   2.798
sequence_randomisationPM  0.63158    0.25267   2.500

Correlation of Fixed Effects:
            (Intr) trtmnM period
trtmntMdfnl -0.130              
periode     -0.389  0.000       
sqnc_rndmPM -0.645  0.000  0.000
\end{verbatim}

\textbf{Interprétation détaillée des outputs du Modèle Mixte (FSS) :} Le
retour d'un objet \texttt{lmerTest} comporte trois noyaux cruciaux : -
\textbf{Random Effects (Variabilité individuelle) :} La valeur de
variance d'\texttt{(Intercept)\ \textbar{}\ id\_patient} prouve
l'existence d'une hétérogénéité intrinsèque large dès l'entrée en essai,
justifiant le recours au modèle mixte. - \textbf{Fixed Effects (Le cœur
de l'analyse) :} - \textbf{\texttt{(Intercept)} (\(4.97\)) :} Le Score
FSS fondamental (Placebo, Période 1). -
\textbf{\texttt{traitementModafinil} (\(-0.989\)) :} C'est le résultat
majeur. Toutes dimensions fixées, le passage au Modafinil octroie une
réduction de pratiquement \textbf{\(-0.99\) point} à l'échelle de
gravité de la fatigue. La \(t-value (-19.48)\) mène impérativement à une
p-value minuscule (\(Pr(>|t|) < 0.001\)). Cette découverte rend
irréfutable un impact curatif statistique de premier plan de l'agent
Modafinil par rapport au Placebo. - \textbf{\texttt{periode}
(\(+0.142\), \(p = 0.008\)) :} Confirme très subtilement un léger
gradient de sur-fatigue chronologique au fil du temps. -
\textbf{\texttt{sequence\_randomisationPM} (\(p = 0.017\)) :} Confirme
la présence d'un ``effet carry-over'', validant son insertion dans le
modèle de correction.

\section{Analyses des critères
secondaires}\label{analyses-des-crituxe8res-secondaires}

\subsection{Évolution de la Visual Analog Scale for Fatigue
(VASF)}\label{uxe9volution-de-la-visual-analog-scale-for-fatigue-vasf}

\begin{Shaded}
\begin{Highlighting}[]
\NormalTok{modele\_vasf }\OtherTok{\textless{}{-}}\NormalTok{ lme4}\SpecialCharTok{::}\FunctionTok{lmer}\NormalTok{(}
\NormalTok{  vasf }\SpecialCharTok{\textasciitilde{}}\NormalTok{ traitement }\SpecialCharTok{+}\NormalTok{ periode }\SpecialCharTok{+}\NormalTok{ visite }\SpecialCharTok{+}\NormalTok{ (}\DecValTok{1} \SpecialCharTok{|}\NormalTok{ id\_patient),}
  \AttributeTok{data =} \FunctionTok{subset}\NormalTok{(df\_long\_simplifie, visite }\SpecialCharTok{\%in\%} \FunctionTok{c}\NormalTok{(}\StringTok{"S3"}\NormalTok{, }\StringTok{"S6"}\NormalTok{, }\StringTok{"S11"}\NormalTok{, }\StringTok{"S14"}\NormalTok{))}
\NormalTok{)}
\FunctionTok{summary}\NormalTok{(modele\_vasf)}
\end{Highlighting}
\end{Shaded}

\begin{verbatim}
Linear mixed model fit by REML ['lmerMod']
Formula: vasf ~ traitement + periode + visite + (1 | id_patient)
   Data: subset(df_long_simplifie, visite %in% c("S3", "S6", "S11", "S14"))

REML criterion at convergence: 1063.9

Scaled residuals: 
     Min       1Q   Median       3Q      Max 
-2.67067 -0.58493  0.07962  0.53566  2.54438 

Random effects:
 Groups     Name        Variance Std.Dev.
 id_patient (Intercept) 114.88   10.718  
 Residual                37.62    6.133  
Number of obs: 152, groups:  id_patient, 38

Fixed effects:
                    Estimate Std. Error t value
(Intercept)         62.91842    3.82823  16.435
traitementModafinil -8.96842    0.99497  -9.014
periode              5.99737    2.22483   2.696
visite.L            -4.78401    2.22483  -2.150
visite.Q            -0.08684    0.99497  -0.087

Correlation of Fixed Effects:
            (Intr) trtmnM period vist.L
trtmntMdfnl -0.130                     
periode     -0.872  0.000              
visite.L     0.780  0.000 -0.894       
visite.Q     0.000  0.000  0.000  0.000
fit warnings:
fixed-effect model matrix is rank deficient so dropping 1 column / coefficient
\end{verbatim}

\textbf{Interprétation détaillée de l'output (VASF) :} - \textbf{Effet
du Traitement :} Analysons la ligne \texttt{traitementModafinil} au sein
du panel \emph{Fixed effects}. Son paramètre \texttt{Estimate} affiche
une baisse de \textbf{\(-8.968\)} avec un test \texttt{t-value}
extrêmement fort (\(-9.01\)). - \textbf{Significativité :} La p-value
associée (\(7.08 \times 10^{-15}\)) valide très rigoureusement
l'hypothèse d'une dissemblance marquée. En résumé, le Modafinil diminue
l'échelle VASF perçue par le patient d'environ \textbf{9 points
complets}, consolidant l'efficacité montrée par le FSS.

\subsection{Évolution de la Fatigue Impact Scale
(FIS)}\label{uxe9volution-de-la-fatigue-impact-scale-fis}

\begin{Shaded}
\begin{Highlighting}[]
\NormalTok{modele\_fis }\OtherTok{\textless{}{-}}\NormalTok{ lme4}\SpecialCharTok{::}\FunctionTok{lmer}\NormalTok{(}
\NormalTok{  fis }\SpecialCharTok{\textasciitilde{}}\NormalTok{ traitement }\SpecialCharTok{+}\NormalTok{ periode }\SpecialCharTok{+}\NormalTok{ visite }\SpecialCharTok{+}\NormalTok{ (}\DecValTok{1} \SpecialCharTok{|}\NormalTok{ id\_patient),}
  \AttributeTok{data =} \FunctionTok{subset}\NormalTok{(df\_long\_simplifie, visite }\SpecialCharTok{\%in\%} \FunctionTok{c}\NormalTok{(}\StringTok{"S3"}\NormalTok{, }\StringTok{"S6"}\NormalTok{, }\StringTok{"S11"}\NormalTok{, }\StringTok{"S14"}\NormalTok{))}
\NormalTok{)}
\FunctionTok{summary}\NormalTok{(modele\_fis)}
\end{Highlighting}
\end{Shaded}

\begin{verbatim}
Linear mixed model fit by REML ['lmerMod']
Formula: fis ~ traitement + periode + visite + (1 | id_patient)
   Data: subset(df_long_simplifie, visite %in% c("S3", "S6", "S11", "S14"))

REML criterion at convergence: 1203.2

Scaled residuals: 
     Min       1Q   Median       3Q      Max 
-2.75774 -0.67395  0.03243  0.58260  2.06385 

Random effects:
 Groups     Name        Variance Std.Dev.
 id_patient (Intercept) 231.9    15.23   
 Residual               104.4    10.22   
Number of obs: 152, groups:  id_patient, 38

Fixed effects:
                    Estimate Std. Error t value
(Intercept)          100.980      6.196  16.297
traitementModafinil  -15.568      1.658  -9.391
periode               11.308      3.707   3.050
visite.L             -11.457      3.707  -3.091
visite.Q              -1.926      1.658  -1.162

Correlation of Fixed Effects:
            (Intr) trtmnM period vist.L
trtmntMdfnl -0.134                     
periode     -0.897  0.000              
visite.L     0.803  0.000 -0.894       
visite.Q     0.000  0.000  0.000  0.000
fit warnings:
fixed-effect model matrix is rank deficient so dropping 1 column / coefficient
\end{verbatim}

\textbf{Interprétation détaillée de l'output (FIS - Impact au quotidien)
:} - \textbf{\texttt{traitementModafinil} :} La limitation incombée à la
fatigue (le score FIS) est formellement amoindrie de \textbf{\(-15.568\)
points} sous Modafinil (\(p < 0.001\)). Réduire l'impact de 15.5 points
dans un questionnaire évaluant les limitations quotidiennes atteste d'un
bénéfice prononcé et très pragmatique pour l'autonomie des patients. -
\textbf{L'effet visite :} Notons aussi la tendance \texttt{visite.L}
(\(-11.45\), \(p=0.002\)) dénotant une baisse générale de la pénibilité
au fil des mesures.

\subsection{Qualité de Vie (SF-36)}\label{qualituxe9-de-vie-sf-36}

\begin{Shaded}
\begin{Highlighting}[]
\NormalTok{modele\_sf36\_pcs }\OtherTok{\textless{}{-}}\NormalTok{ lme4}\SpecialCharTok{::}\FunctionTok{lmer}\NormalTok{(}
\NormalTok{  sf36\_pcs }\SpecialCharTok{\textasciitilde{}}\NormalTok{ traitement }\SpecialCharTok{+}\NormalTok{ periode }\SpecialCharTok{+}\NormalTok{ visite }\SpecialCharTok{+}\NormalTok{ (}\DecValTok{1} \SpecialCharTok{|}\NormalTok{ id\_patient),}
  \AttributeTok{data =} \FunctionTok{subset}\NormalTok{(df\_long\_simplifie, visite }\SpecialCharTok{\%in\%} \FunctionTok{c}\NormalTok{(}\StringTok{"S3"}\NormalTok{, }\StringTok{"S6"}\NormalTok{, }\StringTok{"S11"}\NormalTok{, }\StringTok{"S14"}\NormalTok{))}
\NormalTok{)}
\FunctionTok{summary}\NormalTok{(modele\_sf36\_pcs)}
\end{Highlighting}
\end{Shaded}

\begin{verbatim}
Linear mixed model fit by REML ['lmerMod']
Formula: sf36_pcs ~ traitement + periode + visite + (1 | id_patient)
   Data: subset(df_long_simplifie, visite %in% c("S3", "S6", "S11", "S14"))

REML criterion at convergence: 1037.7

Scaled residuals: 
     Min       1Q   Median       3Q      Max 
-2.59098 -0.61078  0.01787  0.52881  2.06041 

Random effects:
 Groups     Name        Variance Std.Dev.
 id_patient (Intercept) 51.87    7.202   
 Residual               37.61    6.133   
Number of obs: 152, groups:  id_patient, 38

Fixed effects:
                    Estimate Std. Error t value
(Intercept)          48.9408     3.6049  13.576
traitementModafinil   7.6526     0.9949   7.692
periode              -6.4632     2.2247  -2.905
visite.L              6.0315     2.2247   2.711
visite.Q             -0.5053     0.9949  -0.508

Correlation of Fixed Effects:
            (Intr) trtmnM period vist.L
trtmntMdfnl -0.138                     
periode     -0.926  0.000              
visite.L     0.828  0.000 -0.894       
visite.Q     0.000  0.000  0.000  0.000
fit warnings:
fixed-effect model matrix is rank deficient so dropping 1 column / coefficient
\end{verbatim}

\textbf{Interprétation détaillée de l'output (SF-36 Composante Physique)
:} \emph{(À l'inverse des échelles de fatigue, viser une valeur SF-36
grandissante est corrélée à une MEILLEURE fonctionnalité).} -
\textbf{Effet du Traitement :} L'estimation est ici \textbf{POSITIVE
(\(+7.652\))} avec une \texttt{t-value} poussée à \(7.69\) conduisant à
un \(p < 0.001\). - \textbf{Conclusion détaillée :} Sous traitement par
Modafinil, la note de la composante corporelle de qualité de vie remonte
de 7.6 points. Cela prouve que le soulagement induit sur la fatigue se
traduit concrètement par un regain net d'état de santé physique global.

\section{Tolérance au traitement et
Sécurité}\label{toluxe9rance-au-traitement-et-suxe9curituxe9}

Nous étudions l'occurrence d'événements indésirables (modélisés en
régression logistique à effets mixtes).

\begin{Shaded}
\begin{Highlighting}[]
\NormalTok{df\_long\_simplifie}\SpecialCharTok{$}\NormalTok{ei\_num }\OtherTok{\textless{}{-}} \FunctionTok{as.numeric}\NormalTok{(df\_long\_simplifie}\SpecialCharTok{$}\NormalTok{ei)}
\NormalTok{modele\_ei }\OtherTok{\textless{}{-}}\NormalTok{ lme4}\SpecialCharTok{::}\FunctionTok{glmer}\NormalTok{(}
\NormalTok{  ei\_num }\SpecialCharTok{\textasciitilde{}}\NormalTok{ traitement }\SpecialCharTok{+}\NormalTok{ periode }\SpecialCharTok{+}\NormalTok{ visite }\SpecialCharTok{+}\NormalTok{ (}\DecValTok{1} \SpecialCharTok{|}\NormalTok{ id\_patient),}
  \AttributeTok{data =} \FunctionTok{subset}\NormalTok{(df\_long\_simplifie, visite }\SpecialCharTok{\%in\%} \FunctionTok{c}\NormalTok{(}\StringTok{"S3"}\NormalTok{, }\StringTok{"S6"}\NormalTok{, }\StringTok{"S11"}\NormalTok{, }\StringTok{"S14"}\NormalTok{)),}
  \AttributeTok{family =} \FunctionTok{binomial}\NormalTok{()}
\NormalTok{)}
\FunctionTok{summary}\NormalTok{(modele\_ei)}
\end{Highlighting}
\end{Shaded}

\begin{verbatim}
Generalized linear mixed model fit by maximum likelihood (Laplace
  Approximation) [glmerMod]
 Family: binomial  ( logit )
Formula: ei_num ~ traitement + periode + visite + (1 | id_patient)
   Data: subset(df_long_simplifie, visite %in% c("S3", "S6", "S11", "S14"))

      AIC       BIC    logLik -2*log(L)  df.resid 
    137.0     155.1     -62.5     125.0       146 

Scaled residuals: 
    Min      1Q  Median      3Q     Max 
-0.7384 -0.3712 -0.2858 -0.2322  3.8842 

Random effects:
 Groups     Name        Variance Std.Dev.
 id_patient (Intercept) 1.069    1.034   
Number of obs: 152, groups:  id_patient, 38

Fixed effects:
                    Estimate Std. Error z value Pr(>|z|)  
(Intercept)          -3.9109     1.7123  -2.284   0.0224 *
traitementModafinil   0.9374     0.5074   1.848   0.0647 .
periode               0.8892     1.0469   0.849   0.3957  
visite.L             -1.0608     1.1039  -0.961   0.3366  
visite.Q             -0.4744     0.4937  -0.961   0.3366  
---
Signif. codes:  0 '***' 0.001 '**' 0.01 '*' 0.05 '.' 0.1 ' ' 1

Correlation of Fixed Effects:
            (Intr) trtmnM period vist.L
trtmntMdfnl -0.246                     
periode     -0.951  0.037              
visite.L     0.849 -0.042 -0.882       
visite.Q     0.029 -0.042  0.012  0.052
fit warnings:
fixed-effect model matrix is rank deficient so dropping 1 column / coefficient
\end{verbatim}

\textbf{Interprétation détaillée de l'output (Tolérance / EI) Modèle
\texttt{glmer} :} La variable des Évènements Indésirables étant binaire
(Oui=1 / Non=0), l'analyse recourt à un modèle de régression logistique
Mixte (famille \texttt{binomial}). - \textbf{\texttt{Estimate}
(Log-Odds) :} La lecture des effets fixes indique des Log-Odds (des
transformations de rapports de cotes). -
\textbf{\texttt{traitementModafinil} (\(+0.937\)) :} Le modafinil
accroît la variable Log-Odds de \(0.937\). - \textbf{Calcul de l'Odds
Ratio :} La conversion exponentielle donne l'Odds Ratio :
\(OR = \exp(0.937) \approx 2.55\). Sous Modafinil, le risque relatif
d'EI est donc multiplié par environ 2.55. - \textbf{Test de
Significativité :} Toutefois, sa \texttt{p-value}
(\texttt{Pr(\textgreater{}\textbar{}z\textbar{})} = \(0.0647\)) plafonne
\textbf{légèrement au-dessus du seuil strict \(\alpha=0.05\)}. -
\textbf{Conclusion clinique :} Il existe une \emph{tendance} à observer
plus d'événements indésirables avec le Modafinil. Bien que ce ne soit
pas formellement significatif au seuil de 5\%, le signal \(p=0.065\)
inciterait le prescripteur à appliquer une pharmacovigilance attentive
face aux bénéfices physiques constatés.




\end{document}
